% Options for packages loaded elsewhere
\PassOptionsToPackage{unicode}{hyperref}
\PassOptionsToPackage{hyphens}{url}
\documentclass[
]{article}
\usepackage{xcolor}
\usepackage[margin=1in]{geometry}
\usepackage{amsmath,amssymb}
\setcounter{secnumdepth}{-\maxdimen} % remove section numbering
\usepackage{iftex}
\ifPDFTeX
  \usepackage[T1]{fontenc}
  \usepackage[utf8]{inputenc}
  \usepackage{textcomp} % provide euro and other symbols
\else % if luatex or xetex
  \usepackage{unicode-math} % this also loads fontspec
  \defaultfontfeatures{Scale=MatchLowercase}
  \defaultfontfeatures[\rmfamily]{Ligatures=TeX,Scale=1}
\fi
\usepackage{lmodern}
\ifPDFTeX\else
  % xetex/luatex font selection
\fi
% Use upquote if available, for straight quotes in verbatim environments
\IfFileExists{upquote.sty}{\usepackage{upquote}}{}
\IfFileExists{microtype.sty}{% use microtype if available
  \usepackage[]{microtype}
  \UseMicrotypeSet[protrusion]{basicmath} % disable protrusion for tt fonts
}{}
\makeatletter
\@ifundefined{KOMAClassName}{% if non-KOMA class
  \IfFileExists{parskip.sty}{%
    \usepackage{parskip}
  }{% else
    \setlength{\parindent}{0pt}
    \setlength{\parskip}{6pt plus 2pt minus 1pt}}
}{% if KOMA class
  \KOMAoptions{parskip=half}}
\makeatother
\usepackage{color}
\usepackage{fancyvrb}
\newcommand{\VerbBar}{|}
\newcommand{\VERB}{\Verb[commandchars=\\\{\}]}
\DefineVerbatimEnvironment{Highlighting}{Verbatim}{commandchars=\\\{\}}
% Add ',fontsize=\small' for more characters per line
\usepackage{framed}
\definecolor{shadecolor}{RGB}{248,248,248}
\newenvironment{Shaded}{\begin{snugshade}}{\end{snugshade}}
\newcommand{\AlertTok}[1]{\textcolor[rgb]{0.94,0.16,0.16}{#1}}
\newcommand{\AnnotationTok}[1]{\textcolor[rgb]{0.56,0.35,0.01}{\textbf{\textit{#1}}}}
\newcommand{\AttributeTok}[1]{\textcolor[rgb]{0.13,0.29,0.53}{#1}}
\newcommand{\BaseNTok}[1]{\textcolor[rgb]{0.00,0.00,0.81}{#1}}
\newcommand{\BuiltInTok}[1]{#1}
\newcommand{\CharTok}[1]{\textcolor[rgb]{0.31,0.60,0.02}{#1}}
\newcommand{\CommentTok}[1]{\textcolor[rgb]{0.56,0.35,0.01}{\textit{#1}}}
\newcommand{\CommentVarTok}[1]{\textcolor[rgb]{0.56,0.35,0.01}{\textbf{\textit{#1}}}}
\newcommand{\ConstantTok}[1]{\textcolor[rgb]{0.56,0.35,0.01}{#1}}
\newcommand{\ControlFlowTok}[1]{\textcolor[rgb]{0.13,0.29,0.53}{\textbf{#1}}}
\newcommand{\DataTypeTok}[1]{\textcolor[rgb]{0.13,0.29,0.53}{#1}}
\newcommand{\DecValTok}[1]{\textcolor[rgb]{0.00,0.00,0.81}{#1}}
\newcommand{\DocumentationTok}[1]{\textcolor[rgb]{0.56,0.35,0.01}{\textbf{\textit{#1}}}}
\newcommand{\ErrorTok}[1]{\textcolor[rgb]{0.64,0.00,0.00}{\textbf{#1}}}
\newcommand{\ExtensionTok}[1]{#1}
\newcommand{\FloatTok}[1]{\textcolor[rgb]{0.00,0.00,0.81}{#1}}
\newcommand{\FunctionTok}[1]{\textcolor[rgb]{0.13,0.29,0.53}{\textbf{#1}}}
\newcommand{\ImportTok}[1]{#1}
\newcommand{\InformationTok}[1]{\textcolor[rgb]{0.56,0.35,0.01}{\textbf{\textit{#1}}}}
\newcommand{\KeywordTok}[1]{\textcolor[rgb]{0.13,0.29,0.53}{\textbf{#1}}}
\newcommand{\NormalTok}[1]{#1}
\newcommand{\OperatorTok}[1]{\textcolor[rgb]{0.81,0.36,0.00}{\textbf{#1}}}
\newcommand{\OtherTok}[1]{\textcolor[rgb]{0.56,0.35,0.01}{#1}}
\newcommand{\PreprocessorTok}[1]{\textcolor[rgb]{0.56,0.35,0.01}{\textit{#1}}}
\newcommand{\RegionMarkerTok}[1]{#1}
\newcommand{\SpecialCharTok}[1]{\textcolor[rgb]{0.81,0.36,0.00}{\textbf{#1}}}
\newcommand{\SpecialStringTok}[1]{\textcolor[rgb]{0.31,0.60,0.02}{#1}}
\newcommand{\StringTok}[1]{\textcolor[rgb]{0.31,0.60,0.02}{#1}}
\newcommand{\VariableTok}[1]{\textcolor[rgb]{0.00,0.00,0.00}{#1}}
\newcommand{\VerbatimStringTok}[1]{\textcolor[rgb]{0.31,0.60,0.02}{#1}}
\newcommand{\WarningTok}[1]{\textcolor[rgb]{0.56,0.35,0.01}{\textbf{\textit{#1}}}}
\usepackage{graphicx}
\makeatletter
\newsavebox\pandoc@box
\newcommand*\pandocbounded[1]{% scales image to fit in text height/width
  \sbox\pandoc@box{#1}%
  \Gscale@div\@tempa{\textheight}{\dimexpr\ht\pandoc@box+\dp\pandoc@box\relax}%
  \Gscale@div\@tempb{\linewidth}{\wd\pandoc@box}%
  \ifdim\@tempb\p@<\@tempa\p@\let\@tempa\@tempb\fi% select the smaller of both
  \ifdim\@tempa\p@<\p@\scalebox{\@tempa}{\usebox\pandoc@box}%
  \else\usebox{\pandoc@box}%
  \fi%
}
% Set default figure placement to htbp
\def\fps@figure{htbp}
\makeatother
% definitions for citeproc citations
\NewDocumentCommand\citeproctext{}{}
\NewDocumentCommand\citeproc{mm}{%
  \begingroup\def\citeproctext{#2}\cite{#1}\endgroup}
\makeatletter
 % allow citations to break across lines
 \let\@cite@ofmt\@firstofone
 % avoid brackets around text for \cite:
 \def\@biblabel#1{}
 \def\@cite#1#2{{#1\if@tempswa , #2\fi}}
\makeatother
\newlength{\cslhangindent}
\setlength{\cslhangindent}{1.5em}
\newlength{\csllabelwidth}
\setlength{\csllabelwidth}{3em}
\newenvironment{CSLReferences}[2] % #1 hanging-indent, #2 entry-spacing
 {\begin{list}{}{%
  \setlength{\itemindent}{0pt}
  \setlength{\leftmargin}{0pt}
  \setlength{\parsep}{0pt}
  % turn on hanging indent if param 1 is 1
  \ifodd #1
   \setlength{\leftmargin}{\cslhangindent}
   \setlength{\itemindent}{-1\cslhangindent}
  \fi
  % set entry spacing
  \setlength{\itemsep}{#2\baselineskip}}}
 {\end{list}}
\usepackage{calc}
\newcommand{\CSLBlock}[1]{\hfill\break\parbox[t]{\linewidth}{\strut\ignorespaces#1\strut}}
\newcommand{\CSLLeftMargin}[1]{\parbox[t]{\csllabelwidth}{\strut#1\strut}}
\newcommand{\CSLRightInline}[1]{\parbox[t]{\linewidth - \csllabelwidth}{\strut#1\strut}}
\newcommand{\CSLIndent}[1]{\hspace{\cslhangindent}#1}
\setlength{\emergencystretch}{3em} % prevent overfull lines
\providecommand{\tightlist}{%
  \setlength{\itemsep}{0pt}\setlength{\parskip}{0pt}}
\usepackage{bookmark}
\IfFileExists{xurl.sty}{\usepackage{xurl}}{} % add URL line breaks if available
\urlstyle{same}
\hypersetup{
  pdftitle={draft\_manuscript\_Dual\_Coil},
  hidelinks,
  pdfcreator={LaTeX via pandoc}}

\title{draft\_manuscript\_Dual\_Coil}
\author{}
\date{\vspace{-2.5em}2026-02-12}

\begin{document}
\maketitle

\begin{Shaded}
\begin{Highlighting}[]
\FunctionTok{list.files}\NormalTok{(}\FunctionTok{getwd}\NormalTok{(), }\AttributeTok{pattern =} \StringTok{".html"}\NormalTok{)}
\end{Highlighting}
\end{Shaded}

\begin{verbatim}
## [1] "draft_manuscript.html"
\end{verbatim}

\subsection{Introduction}\label{introduction}

\subsubsection{Sensory Integration and Motor Planning (Sample Intro
First
Paragraph)}\label{sensory-integration-and-motor-planning-sample-intro-first-paragraph}

Activity in the posterior parietal cortex (PPC) is associated with key
roles in motor planning prior to movement execution, integrating visual,
auditory, and somatosensory information prior to movement initiation.
Activity in the PPC is associated with the initial planning of a
physical action (Andersen \& Buneo, 2002; Cohen \& Andersen, 2002){]}.
During the earliest stages of motor planning prior to execution, the PPC
is involved with the transformation and coordination of reference-frames
from various sensory modalities; integrating visual coordinates
originating from the eyes, auditory information related to head
coordinates, and tactile coordinates from the body. Neurons located in
the PPC are associated with the transformation of these varied sensory
stimuli into general coordinates that form a common frame of reference,
essential for the coordination of movements of different effectors to
ensure coherent movement execution (Andersen \& Buneo, 2002). Early
stages of motor planning taking place in the PPC represent the expected
visual outcomes of a movement prior to downstream motor execution, with
later stages of motor activity occurring primarily in the primary motor
cortex (M1) (Pilacinski et al., 2018). There is evidence that suggests
the PPC, independent of the M1, processes visual information used to
later produce a corresponding motor output (Pilacinski et al., 2018),
further implicating the PPC's role in the visual-integration aspects of
motor planning/preparation.

\subsubsection{Research Aim}\label{research-aim}

To determine the intrahemispheric interconnectivity of the PPC
(specifically the supramarginal gyrus) with the M1 during motor imagery
(MI; the imagination of movement without physical execution) using
dual-coil TMS.

\subsection{Methods (Sample)}\label{methods-sample}

Participants completed blocks of MI of a finger abduction task while
dual-coil TMS was administered over the left PPC and M1.Electromyography
(EMG) of the finger muscle was used to assess corticospinal excitability
during the task.Paired TMS pulses were administered with the PPC as
priming stimuli, followed by the test stimuli to the M1 after 4/8
milliseconds interstimulus interval (ISI).Single pulses were
administered for comparison to the paired pulses.

\subsubsection{Data Analysis}\label{data-analysis}

Peak-to-peak amplitudes of motor evoked potentials (MEPs) elicited by
TMS during the trials were analyzed to determine facilitatory or
inhibitory interactions between the PPC and the M1 during MI.Comparisons
were made between paired-TMS trials (PPC and M1 pulses) and M1 stimulus
only trials.

\subsection{Figures/Results Draft}\label{figuresresults-draft}

Overall, group level analyses yielded no significant interactions
between the PPC and M1 during MI using our methodology. However,
individual participant data showed individual variation in
parietal-motor interaction during MI potential providing evidence for
modulating activity at the 4ms ISI in differing ways for different
individuals.

\begin{figure}
\centering
\pandocbounded{\includegraphics[keepaspectratio]{draft_manuscript---Copy_files/figure-latex/unnamed-chunk-3-1.pdf}}
\caption{MEP amplitude normalised to M1-only mean}
\end{figure}

\subsection{Discussion}\label{discussion}

Findings provide insight into the neurophysiology underlying MI as well
as dual-coil TMS stimulation parameters in research.

\subsubsection{Limitations}\label{limitations}

Due to the nature of the setup, we were unable to hotspot the PPC for
each individual participant to guarantee we stimulate the same place on
their brain. This could Potentially account for individual variability
due to slight differences in TMS stimulus location.We are also limited
to analysing a small sample size at the present moment.

\section*{References}\label{references}
\addcontentsline{toc}{section}{References}

\phantomsection\label{refs}
\begin{CSLReferences}{1}{0}
\bibitem[\citeproctext]{ref-andersenIntentionalMapsPosterior2002}
Andersen, R. A., \& Buneo, C. A. (2002). Intentional maps in posterior
parietal cortex. \emph{Annual Review of Neuroscience}, \emph{25},
189--220. \url{https://doi.org/10.1146/annurev.neuro.25.112701.142922}

\bibitem[\citeproctext]{ref-cohenCommonReferenceFrame2002}
Cohen, Y. E., \& Andersen, R. A. (2002). A common reference frame for
movement plans in the posterior parietal cortex. \emph{Nature Reviews.
Neuroscience}, \emph{3}(7), 553--562.
\url{https://doi.org/10.1038/nrn873}

\bibitem[\citeproctext]{ref-pilacinskiHumanPosteriorParietal2018}
Pilacinski, A., Wallscheid, ⨯. M., \& Lindner, A. (2018). Human
posterior parietal and dorsal premotor cortex encode the visual
properties of an upcoming action. \emph{PLoS One}, \emph{13}(10),
e0198051. \url{https://doi.org/10.1371/journal.pone.0198051}

\end{CSLReferences}

\end{document}
